\documentclass[../../main.tex]{subfiles}
\begin{document}

{\bf [解]}

題意の“ある時刻”に自動車が原点にいるとし，自動車は $x$ 軸正方向へ動くものとする．また飛行機は平面 $z=k\ \text{[km]}\ (>0)$ 上を飛ぶとする．
$3$つの点を
\begin{align}
  \begin{cases}
    A \cdots \text{“ある時刻”における飛行機の座標．} yz \text{平面上にとる．} \\
    B \cdots 36 \text{秒後における飛行機の座標．}  \\
    C \cdots 36 \text{秒後における自動車の座標}
  \end{cases}
\end{align}
とする．また点 $X$ に対し，$X$ の $xy$ 平面への正射影を $X'$ で表す．
$B'$から$x$軸におろした垂足を$B''$とする．
図は以下のようになる．

\begin{figure}[htb]
  \centering
  \begin{tikzpicture}[
      scale=1.4,
      >=stealth,
      every node/.style={font=\small},
      x={(0.9cm, -0.25cm)},   % x軸（東・車の進行方向）
      y={(0.55cm, 0.45cm)},   % y軸（北方向）
      z={(0cm, 1.35cm)}       % z軸（鉛直上向き・高度）
    ]

    % --- 1. 数値・パラメータ設定 (k = 1.0) ---
    % A(0, sqrt(3), 1)
    % C(1, 0, 0)
    % B(1 + 1.5, sqrt(3)/2, 1) = (2.5, 0.866, 1)
    \def\kVal{1.0}
    \pgfmathsetmacro{\sqrtThree}{sqrt(3)}

    % 各点の空間座標
    \coordinate (O)   at (0, 0, 0);
    \coordinate (Ap)  at (0, \sqrtThree, 0);
    \coordinate (A)   at (0, \sqrtThree, \kVal);
    \coordinate (C)   at (1.0, 0, 0);
    \coordinate (Bp)  at (2.5, {\sqrtThree/2}, 0);
    \coordinate (B)   at (2.5, {\sqrtThree/2}, \kVal);
    \coordinate (Bpp) at (2.5,0,0);

    % --- 2. 座標軸 ---
    \draw[->, thick] (-0.3, 0, 0) -- (3.3, 0, 0) node[below right] {$x$ (北)};
    \draw[->, thick] (0, -0.3, 0) -- (0, 2.5, 0) node[above right] {$y$ (西)};
    \draw[->, thick] (0, 0, -0.2) -- (0, 0, 1.7) node[above] {$z$ (上)};
    \node[below left=2pt] at (O) {$O$};

    % --- 3. xy 平面上の補助線（破線） ---
    \draw[thick, dashed, gray!70] (O) -- (Ap) node[midway, left=2pt] {$\sqrt{3}k$};
    \draw[thick, dashed, gray!70] (C) -- (Bp);
    \draw[dashed, gray!50] (Ap) -- (Bp);
    \draw[thick,dashed,gray!50] (Bp) -- (Bpp);

    % C から北向き（y軸平行）の基準点線（方位角表示用）
    \draw[dashed, gray!60] (C) -- ++(0, 1.2, 0);

    % --- 4. 鉛直垂線（高度 k） ---
    \draw[thick, dashed, blue!70!black] (Ap) -- (A) node[midway, left=1pt] {$k$};
    \draw[thick, dashed, blue!70!black] (Bp) -- (B) node[midway, right=1pt] {$k$};

    % --- 5. 視線および移動軌跡（実線） ---
    % t = 0 での視線 OA
    \draw[very thick, orange!80!black] (O) -- (A);
    % 36秒後での視線 CB
    \draw[very thick, orange!80!black] (C) -- (B);
    % 飛行機の軌跡 AB
    \draw[very thick, magenta!90!black] (A) -- (B)
    node[midway, above, font=\footnotesize] {$|AB|=\sqrt{7}$};
    % 自動車の軌跡 OC
    \draw[very thick, cyan!80!black] (O) -- (C)
    node[midway, below, font=\footnotesize] {$1\text{km}$};

    % --- 6. 直角記号 ---
    % A' における垂直
    \draw[gray!70] ($ (Ap) + (0, -0.15, 0) $) -- ($ (Ap) + (0, -0.15, 0.15) $) -- ($ (Ap) + (0, 0, 0.15) $);
    % B' における垂直
    \draw[gray!70] ($ (Bp) + (-0.15, 0, 0) $) -- ($ (Bp) + (-0.15, 0, 0.15) $) -- ($ (Bp) + (0, 0, 0.15) $);

    % --- 7. 角度の円弧とラベル ---
    % (1) t = 0 での仰角 \angle AOA' = 30°
    \draw[thick, green!60!black] ($ (O)!0.32!(Ap) $) to[bend right=15] ($ (O)!0.32!(A) $);
    \node[green!60!black, font=\footnotesize, left=2pt] at ($ (O)!0.35!(A) + (0, 0.05, 0.05) $) {$30^\circ$};

    % (2) 36秒後 C での仰角 \angle BCB' = 30°
    \draw[thick, green!60!black] ($ (C)!0.28!(Bp) $) to[bend left=-18] ($ (C)!0.28!(B) $);
    \node[green!60!black, font=\footnotesize, right=1pt] at ($ (C)!0.3!(B) $) {$30^\circ$};

    % (3) 36秒後 C での仰角 \angle B'CB'' = 30°
    \draw[thick, green!60!black] ($ (C)!0.28!(Bpp) $) to[bend left=-18] ($ (C)!0.28!(Bp) $);
    \node[green!60!black, font=\footnotesize, right=1pt] at ($ (C)!0.3!(Bpp) $) {$30^\circ$};

    % --- 8. 点のプロットとラベル ---
    \fill (O) circle (1.2pt);
    \fill (Ap) circle (1.2pt) node[below=2pt] {$A'$};
    \fill[blue!80!black] (A) circle (1.4pt) node[above left] {$A(0, \sqrt{3}k, k)$};
    \fill[cyan!80!black] (C) circle (1.4pt) node[below=2pt] {$C(1,0,0)$};
    \fill (Bp) circle (1.2pt) node[below right] {$B'$};
    \fill[blue!80!black] (B) circle (1.4pt) node[above right] {$B$};
    \fill[blue!80!black] (Bpp) circle (1.4pt) node[below right] {$B''$};

  \end{tikzpicture}
  \caption{自動車（$O \to C$）と飛行機（$A \to B$）の空間配置および仰角関係}
  \label{fig:airplane_car_3d}
\end{figure}

題意より，$t=0$での車から飛行機への仰角，すなわち $\angle A O A'$が $30^\circ$ だから $A(0, \sqrt{3}k, k)$ である．車と飛行機の速度を秒速に直すと，
\begin{align}
  100\ \text{[km/h]} = \dfrac{10^3}{36}\ \text{[m/s]} \\
  \sqrt{7} \times 100\ \text{[km/h]} = \dfrac{\sqrt{7} \cdot 10^3}{36}\ \text{[m/s]}
\end{align}
だから，$C(1 \text{[km]},0,0)$である．

この時点での車から飛行機の角度は北から30度西の方向に仰角$30^{\circ}$である．
従って，まず
\begin{align}
  |B'C| = \dfrac{|BB'|}{\tan \angle BCB'} = \dfrac{k}{\tan 30^{\circ}} = \sqrt{3}k \label{eq:1}
\end{align}
である．よって
\begin{align}
  |B''C| = |B'C|\cos\angle B'CB'' = \cos\angle 30^{\circ} \sqrt{3}k = \dfrac{3}{2}k
\end{align}
および
\begin{align}
  |B'B''| = |B'C|\sin\angle B'CB'' = \sin\angle 30^{\circ} \sqrt{3}k = \dfrac{\sqrt{3}}{2}k
\end{align}
だから，$B$の座標は
\begin{align}
  B\Bigl( 1 + \dfrac{3}{2}k, \dfrac{\sqrt{3}}{2}k, k \Bigr)
\end{align}
である．したがって，
\begin{align}
  \vec{AB} =
  \begin{pmatrix} 1 + \dfrac{3}{2}k \\ -\dfrac{\sqrt{3}}{2}k \\ 0
  \end{pmatrix}
\end{align}
となる．題意より$|\vec{AB}| = \sqrt{7}\ \text{[km]}$ だから，2乗して整理して，
\begin{align}
  7 &= |AB|^2 \\
  &= \Bigl(1 + \dfrac{3}{2}k\Bigr)^2 + \Bigl(\dfrac{\sqrt{3}}{2}k\Bigr)^2 \\
  &= 3k^2 + 3k + 1 \\
  \therefore
  k &= 1\ (>0)\ \text{[km]}
\end{align}
である．したがって，求める高度は $k = 1000\ \text{[m]}$ である．

\newpage

\end{document}
